This report presents work carried out within the CareBot projects at HAN
University of Applied Sciences in collaboration with
\href{https://ambyon.care}{Ambyon}. The project addresses the need for reliable
robot perception of elevator buttons in hospital environments, where autonomous
navigation across floors depends on accurate detection, interpretation, and
spatial localization of panel interfaces.

The developed workflow combines instance segmentation for button localization,
semi-automatic annotation with iterative retraining, and a lightweight optical
character recognition model for semantic button interpretation. The original
elevator-panel dataset was extended with instance-level annotations and later
augmented with deployment-specific images from the target university
environment. Detector and recognizer models were trained, tuned, and evaluated
using held-out test data and deployment-oriented runtime benchmarks.

The results show that the final system can combine button segmentation, text
recognition, and depth-based target estimation in a modular embedded perception
pipeline running on NVIDIA Jetson hardware. The project therefore delivers a
practical foundation for elevator-button interaction on the CareBot platform and
for future model improvement through deployment-time data collection and
retraining.
